\uuid{trjv}
\exo7id{181}
\titre{exo7 181}
\auteur{gourio}
\organisation{exo7}
\datecreate{2001-09-01}
\isIndication{false}
\isCorrection{false}
\chapitre{Logique, ensemble, raisonnement}
\sousChapitre{Autre}
\cpge{mpsi-mp2i.raisonnement-vocabulaire-ensembliste.a, pcsi.raisonnement-vocabulaire-ensembliste.a, ptsi.raisonnement-vocabulaire-ensembliste.a, mpsi-mp2i.calcul-algebrique-trigonometrie.a, pcsi.calcul-algebrique-trigonometrie.a, ptsi.calcul-algebrique-trigonometrie.a, mpsi-mp2i.arithmetique-entiers.d}
\module{Algèbre}
\niveau{L1}
\difficulte{}

\contenu{
\question{
Montrer que :
$$\forall n\geq 2,u_{n}=\sum\limits_{k=1}^{n}\frac{1}{k}\notin \Nn.$$
\emph{Indication} : montrer que
$$\forall n\geq 2,\exists (p_{n},q_{n})\in (\Nn^{*})^{2},u_{n}=\frac{2p_{n}+1}{2q_{n}}.$$
}
}
